\section{Introduction}

Vision-language models (VLMs) are increasingly used in settings where images are accompanied by textual context: document images with OCR, screenshots with metadata, retrieved descriptions, captions, or user-provided hints.
Such text can be helpful when it is correct, but it can also become a shortcut.
When textual context conflicts with the image, a VLM may answer from the text rather than from visual evidence.
Deng et al.~\cite{deng2025words} describe this phenomenon as ``blind faith in text'', showing that several VLMs disproportionately trust textual information under conflict.

This failure is not merely a benchmark artifact.
In practical systems, auxiliary text is often produced by retrieval, OCR, captioning, or prior model predictions.
All of these sources can be noisy.
If a VLM over-trusts such text, the model can produce confident but visually ungrounded answers.
For document VQA, this is particularly important: the model must decide whether the image or the accompanying text should be trusted when they disagree.

A natural objection is that the model could simply ignore the auxiliary text whenever corruption is possible.
This is not a satisfactory solution.
First, in real use the system is not told whether the text is corrupted, matched, or irrelevant.
Second, auxiliary text is often genuinely useful: when it matches the image, ignoring it can remove information that helps the model answer.
Third, blind text removal does not teach conflict resolution; it only changes the input distribution.
The desired behavior is conditional trust: use text when it is supported by the image, but reject it when it contradicts visual evidence.

We ask a targeted question:
\emph{Can we reduce text-following failures by directly training a VLM to prefer image-grounded answers over its own corrupted-text errors?}
Our key observation is that the corrupted-text setting naturally produces preference pairs.
If a model answers incorrectly and its prediction appears in the corrupted auxiliary text, then the model has likely followed the misleading text.
The ground-truth answer can be used as the preferred response, and the model's text-following prediction can be used as the dispreferred response.
This gives a high-precision preference signal without manual labeling.

We instantiate this idea with Direct Preference Optimization (DPO)~\cite{rafailov2023dpo}.
Unlike online reinforcement learning methods such as PPO or GRPO~\cite{shao2024deepseekmath}, DPO directly optimizes pairwise preferences.
This is well aligned with our setting: the contrastive pair is already available as \{image-grounded answer, text-following wrong answer\}.
Prior work already showed that supervised fine-tuning (SFT) with text augmentation can mitigate blind faith in text~\cite{deng2025words}.
Our goal is not to deny that result, but to ask whether a more targeted preference signal can address the same failure mode.
Instead of training on a broad mixture of text-only, original VQA, match, corruption, and irrelevance samples, we mine the specific cases where the model follows corrupted text and directly contrast that answer with the ground truth.
We compare DPO against GRPO-style sampled-answer optimization, auxiliary SFT preservation, model transfer, and dataset transfer.

Our study is intentionally scoped as an empirical mitigation paper rather than a new benchmark proposal.
Words or Vision already demonstrates that VLMs can over-trust text under controlled conflict.
The remaining question is what kind of training signal most directly changes this behavior.
This distinction matters for the scientific claim.
If the only intervention is to remove text, the result says little about whether a VLM can learn to adjudicate between modalities.
If the intervention is generic SFT, the model may improve because it sees more correct answers, but the mechanism of text-bias reduction remains less explicit.
If the intervention is online RL with a scalar reward, the method depends heavily on candidate sampling and reward shaping.
We therefore focus on preference pairs derived from actual text-following failures, where the chosen and rejected answers instantiate the desired behavioral reversal.

The practical motivation is also broader than document VQA.
Modern VLM pipelines often combine images with OCR outputs, generated captions, retrieved passages, or user-provided textual descriptions.
These textual channels are useful because they can expose information that the model might otherwise miss.
At the same time, they create a failure mode in which a visually plausible but wrong text snippet overrides the image.
A robust VLM should not be text-blind; it should be text-skeptical when the text is unsupported.
This paper evaluates whether DPO can teach this conditional skepticism from a small number of model-specific failures.

Our study makes three empirical findings.
First, on DocVQA with Qwen3-VL-8B, DPO sharply improves corrupted-text robustness while reducing the rate of text-copying among remaining errors.
Second, broadening DPO from corrupted-only text-following errors to multi-condition error-mined pairs yields the strongest result, suggesting that condition-specific error mining can improve stability.
Third, the method has an important boundary condition: it works best when the base VLM already has enough visual/OCR competence.
Qwen-family models benefit substantially, while LLaVA-Next models under our prompt and template show weak DocVQA base performance and limited benefit.

We summarize our contributions as follows:
\begin{itemize}[leftmargin=*]
    \item We define an automatic text-following error mining procedure for corrupted-text VQA, requiring both incorrectness and overlap between the wrong prediction and corrupted auxiliary text.
    \item We show that DPO on these mined pairs substantially reduces text bias on DocVQA and transfers directionally to VQAv2 and Qwen2-VL.
    \item We position DPO relative to the SFT mitigation studied by Deng et al.~\cite{deng2025words}, and compare it with GRPO, auxiliary SFT preservation, multi-condition error mining, and model transfer.
\end{itemize}
